\documentclass[11pt]{article}

\usepackage[margin=1.5in]{geometry}
\usepackage{lecnotes}
% \usepackage{mpass}
% \lstset{style=verb,language=mpass}
\input{lmacros}

\newcommand{\course}{15-316: Software Foundations of Security \& Privacy}
\newcommand{\lecturer}{Matt Fredrikson\footnote{with edits by Giselle Reis, Ryan Riley, and Frank Pfenning}}
\newcommand{\lecdate}{April 16, 2026}
\newcommand{\lecnum}{22}
\newcommand{\lectitle}{Differential Privacy}
\newcommand{\courseurl}{https://15316-cmu.github.io/2026/}

\newcommand{\pr}{\mathrm{Pr}}
\newcommand{\flip}{\mathbf{flip}}

\begin{document}

\maketitle


\section{Introduction}

In \href{\courseurl/lectures/13-declassification.pdf}{Lecture 13} we looked at
ways of relaxing noninterference so that we could reason about the information
flow security of programs that leak some, but not ``too much'' information about
their secrets. To address this problem in the special context of authorization
checks, we discussed a type system due to \citet{Volpano00popl}. They introduced
a $\mathbf{match}$ construct and a corresponding typing rule that we later
generalized to an arbitrary declassification construct.
\begin{tabbing}
\quad $\mb{if}\; \mb{declassify}(\mi{guess} = \mi{pin})\;$ \\
\quad $\mb{then}\; \mi{auth} := 1\;$ \\
\quad $\mb{else}\; \mi{auth} := 0$
\end{tabbing}
For this program, our security policy assigned $\mi{pin}$ high security level
and $\mi{guess}$ and $\mi{auth}$ low security level.

We did not formally introduce the notion then, but we can define the
\emph{feasible set} as those initial states that lead to the observable outcome.
\[
  \Omega_{\Sigma}(\omega, \alpha) =
  \{ \omega' \mid \Sigma \vdash \omega \approx_{\ms{L}} \omega'\ \mbox{and}\
  \Sigma \vdash \eval\; \omega\; \alpha \approx_{\ms{L}} \eval\; \omega'\; \alpha \}
\]
The password checker above is acceptable because its feasible set is rather
large, say, $2^{64} - 1$ if the outcome is $\mi{auth} = 0$.  Every time we run
this program with a new guess we can reduce the size of the feasible set by 1
(even if this temporal evolution isn't part of the formal definition).  By
contrast, if we replaced equality by $\mb{declassify}(\mi{guess} \leq \mi{pin})$
then we can cut the size of the feasible set in half every time we run the
program.  As you saw in \href{\courseurl/labs/lab2-infoflow.pdf}{Lab 2} this
can easily be used for an attack on a server to discover the $\mi{pin}$.

You can think of the size of feasible set as a quantitative measure of the
uncertainty that a user (possibly an attacker) has regarding the high security
part of the initial state, once they know the outcome.

In today's lecture, we will look more carefully at a different set of techniques
for revealing some useful information about secret state while controlling the
attacker's level of uncertainty about it. These techniques all use randomness to
produce approximate results for computations, while providing some form of cover
for the true secret. We will look at a property called \emph{differential
  privacy}~\citep{Dwork06icalp} that formalizes the protections one might gain
from this approach, and study some properties that make it useful for building
computations that protect secret data. Differential privacy has been applied to
a wide range of important computations to protect the privacy of source
data~\citep{Dwork14fttcs}, from machine learning~\citep{Chaudhuri11jmlr} to web browser
data collection~\citep{Erlingsson14ccs}. We will not have time to cover these
applications in any detail, but will instead focus on the core ideas behind the
approach.

\section{Quantifying Uncertainty}

As discussed above, when the feasible set is large, it is an indication that the
associated program did not reveal ``too much'' about its secret initial
state. The reason for this is that when a large number of initial states remain
consistent with an attacker's observations, then the attacker's uncertainty
about which one was actually used is great. So perhaps we can reason about
information flow security in terms of keeping the attacker's uncertainty about
the secret high.

But what can we do if the program that we want to write is inherently ``leaky''
in that it results in small feasible sets? One way that we can make the attacker
more uncertain about the secret initial state is to use randomness in our
program. Consider for example a technique called \emph{randomized
  response}~\citep{Warner65jasa}, which is a privacy technique dating back to the
1960s with roots in the social sciences. Randomized response was motivated by
survey collection, in situations where questions asked of respondents relate to
sensitive issues. Randomized response gives these subjects \emph{plausible
  deniability}, by providing a structured way of adding random ``noise'' to
their answer.

In the following, assume that $\flip()$ is a random function that flips
an unbiased coin. In other words,
\[
\flip() =
\left\{
\begin{array}{ll}
%
1 & \mbox{with probability}\ 1/2 \\
0 & \mbox{with probability}\ 1/2
\end{array}
\right.
\]
Then suppose that $F$ is a function that returns a value in $\{0,1\}$,
and that we wish to release $F(x)$ publicly while hiding the secret
value $x$ as much as possible. Then the randomized response program
$\mathsf{RandResp}$, is as follows, where we assume that the variable
$\mi{out}$ is publicly-observable and $b$ is not (e.g., $\Gamma = x : \ms{H}, b
: \ms{H}, \mi{out} : \ms{L})$.
\[
\begin{array}{l}
%
b := \flip() \\
\mathbf{if}\ b = 1\ \mathbf{then} \\
\ \ \mi{out} := F(x) \\
\mathbf{else} \\
\ \ \mi{out} := \flip()
\end{array}
\]
In short, randomized response returns the true value of $F(x)$ with
probability 1/2, and a completely random answer with probability 1/2. In
terms of feasible sets, this appears to be an absolutely brilliant
approach because now the attacker must be completely uncertain about the
initial value of $x$. Why is this so? The adversary can only see $\mi{out}$,
and if $b = 0$ after being assigned, then $\mi{out}$ does not depend at all on
$x$, so $x$ could be anything as though the program satisfied
non-interference.

But perhaps this does not seem quite right. Let us assume for a moment
that $x \in \{0,1\}$ and $F$ is simply the identity function, and walk
through the various possibilities. In the following, we will treat
$\mathsf{RandResp}$ as though it were a function of $x$ that returns the
value in $\mi{out}$ after executing. If $x = 0$, then,
\[
\pr[\mathsf{RandResp}(0) = 0] = \pr[b = 1] + \pr[b = 0 \land \flip() = 0] = 1/2 + 1/4 = 3/4
\]
We could use the exact same reasoning to conclude that
$\pr[\mathsf{RandResp}(1) = 1] = 3/4$. Likewise we could reason about
the probability that randomized response outputs an incorrect answer,
\[
\pr[\mathsf{RandResp}(0) = 1] = 1 - \pr[\mathsf{RandResp}(0) = 0] = \pr[b = 0 \land \flip() = 1] = 1/4
\]
So we see that $\mathsf{RandResp}$ outputs the \emph{correct} value of
$F(x)$ with fairly high probability of 3/4, and an incorrect ``random''
value with probability 1/4. In other words, most of the time the
attacker is safe in assuming that $\mathsf{RandResp}$ outputs exactly
the same value as $F(x)$, and so can go about inferring $x$ by computing
feasible sets as before.

This is not to say that randomized response does nothing to protect $x$, and
indeed it may offer ample protection for many applications because the attacker
still has more uncertainty than they would otherwise. But by reasoning about the
probabilities of various outcomes and what the attacker is able to infer from
them, we arrived at a much more nuanced view of the degree of security than was
suggested by looking at the feasible set of $\mathsf{RandResp}$ alone.

\subsection{Quantifying a Tradeoff}

There are some arbitrary choices that have been made in this conception
of randomized response, and they influence the degree of adversarial
uncertainty of the secret input $x$. In particular, we could generalize
$\flip()$ by adding a parameter $0 \le p \le 1$ controlling the bias of
the coin.
\[
\flip(p) =
\left\{
\begin{array}{ll}
%
1 & \mbox{with probability}\ p \\
0 & \mbox{with probability}\ 1 - p
\end{array}
\right.
\]
We could use this in $\mathsf{RandResp}$ as follows, assuming $p$ is
chosen to be some constant in advance.
\[
\begin{array}{l}
%
b := \flip(p) \\
\mathbf{if}\ b = 1\ \mathbf{then} \\
\ \ \mi{out} := F(x) \\
\mathbf{else} \\
\ \ \mi{out} := \flip(p)
\end{array}
\]
Then updating the analysis we did before with this more general
solution, we see that:
\[
\begin{array}{l}
  \pr[\mathsf{RandResp}(x) = F(x)] = \pr[b = 1] + \pr[b = 0 \land \flip(p) = F(x)] \\
  \null = p + (1-p)\pr[F(x)=\flip(p)]
\end{array}
\]
When $F$ is the identity function then we have,
\[
  \begin{array}{l}
\pr[\mathsf{RandResp}(0) = 0] = \pr[b = 1] + \pr[b = 0 \land \flip(p) = 0] = p + (1-p)^2 \\
\pr[\mathsf{RandResp}(1) = 1] = \pr[b = 1] + \pr[b = 0 \land \flip(p) = 1] = p + (1-p)p
  \end{array}
\]
So if we set $p \ge 1/2$, then we would be sure to have a more accurate
answer in the sense that $\mathsf{RandResp}$ returns $F(x)$ with greater
likelihood. But this comes at a tradeoff in information flow security,
as the attacker can also be more confident (less uncertain) about the
feasible set. Likewise, smaller values of $p$ lead to a less accurate
solution, but increase the attacker's uncertainty and so afford greater
security.

\section{Differential Privacy}

Now we will turn to a property that is useful in many cases for
characterizing the adversarial uncertainty one obtains through the use
of randomized computation. In this setting, we will assume that the
program $\alpha$ makes use of memory operations, and wants to prevent
too much information about the contents of any cell in from
leaking through its output result. It is called \emph{differential
  privacy}, and is an active area of study and application.

In the following, we will assume that all of the indices in memory $m$
are secret and so typed $\ms{H}$, and that all of the variables used by
the program are typed $\ms{L}$. So intuitively, think of the memory
$m$ as perhaps being an input where each cell holds the data of
one individual that is to be used by $\alpha$. The developer of $\alpha$
wishes to compute some useful aggregate fact about the individuals'
data, and will store the result in the variables of the final state
$\eval\; (\omega, m)\; \alpha$. The goal is to make sure that the
results do not reveal too much information about any single individual's
data stored in $m$.

\begin{definition}[$\epsilon$-Differential Privacy]
\label{def:dp1}
Let $\epsilon \geq 0$. A program $\alpha$ satisfies $\epsilon$-differential
privacy if for all possible memory configurations $m_1$ and
$m_2$ that differ in \emph{exactly one index}, and all states
$\omega$ and $\nu$, the following inequality holds:
\[
\mathrm{Pr}[\eval\; (\omega,m_1)\; \alpha= \nu] \le e^\epsilon \cdot \mathrm{Pr}[\eval\; (\omega,m_2)\; \alpha = \nu]
\]
The probabilities in this expression are taken over the randomness of
$\alpha$'s computation.
\end{definition}

We write $m_1 \sim_1 m_2$ if $m_1$ and $m_2$ differ in one index.  The fact that
$\alpha$ is a program that does not explicitly ``output'' a single value is
indeed irrelevant to the essence of this definition. It may be clearer for some
to just think of $\alpha$ as a function $f$ that takes a memory configuration
$m$ as input and returns a single discrete value rather than a state. This leads
to the following equivalent definition.

\begin{definition}[$\epsilon$-Differential Privacy (functional form)]
\label{def:dp}
Let $\epsilon \geq 0$. A function $f$ satisfies $\epsilon$-differential
privacy if for all possible inputs $m_1$ and $m_2$ with $m_1 \sim_1 m_2$
and all return values $s$ the following inequality holds:
\[
\mathrm{Pr}[f(m_1) = s] \le e^\epsilon \cdot \mathrm{Pr}[f(m_2) = s]
\]
The probabilities in this expression are taken over the randomness of $f$'s
outputs.
\end{definition}

To keep notation as simple as possible, we will stick with the latter
form of the definition for the remainder of the lecture.

First, notice that \autoref{def:dp} is a property of the function
$f$, and \emph{not} of the data being computed on or any particular
output of $f$. In other words, when we speak of something as being
differentially private, we are always referring to a process used to
compute outputs from secret inputs. You may at times hear people refer
to a piece of data as ``differentially private'', but do not get
confused; when used correctly, this language means that the data was
computed by a function that satisfies $\epsilon$-differential privacy.

Second, the $\epsilon$ in \autoref{def:dp} is called the
\emph{privacy budget}, and controls the tradeoff between privacy and
accuracy in much the same way that $p$ did in our randomized response
example before. We will get into some high-level intuitive interpretations
of this definition in a little while, but first let us think about its
various components and how they relate to $f$'s behavior directly.

\paragraph{Privacy budget $\epsilon$.}
We hinted earlier that the privacy budget has an influence on both the
degree of privacy established by the function, as well as the degree of
approximation in the results. $\epsilon$ is our privacy budget, and it
is a numeric real-valued quantity. To understand what it means, let us
look at the behavior of an $\epsilon$-differentially private $f$ for
extremal values of $\epsilon$.

Suppose that we make $\epsilon = 0$. Then \autoref{def:dp}
requires that for any $m_1 \sim_1 m_2$ and outputs $s$,
the stated inequality holds.
Notice that the definition is
symmetric in the values that $m_1$ and $m_2$ take; there is nothing that
distinguishes them from each other, so $f$ must also satisfy:
\[
  \pr[f(m_2) = s] \le \pr[f(m_1) = s]
\]
Combining the two inequalities, it must be that
$\pr[f(m_1) = s] = \pr[f(m_2) = s]$ for all $m_1 \sim_1 m_2$ What does this mean
for the privacy of individuals in $m_1$ and $m_2$, and the utility of $f$?
\begin{itemize}
\item When it comes to privacy, we can conclude that $\epsilon=0$ implies
  \textbf{no leakage} of information about the contents of \emph{any}
  individual. Why does this hold for any index? Recall that \autoref{def:dp}
  needs to hold for \emph{all} pairs $m_1 \sim_1 m_2$. So, if our actual input
  is $m_1$, then all inputs $m_2$ that we obtain by changing one index in $m_1$
  must produce the same distribution of outputs in $f$.

\item As for utility, you probably guessed that $\epsilon = 0$ is not great. In
  fact, because $f$'s output distribution needs to remain the same for all
  adjacent inputs, we can observe that by transitivity $f$'s output distribution
  needs to remain the same for \emph{all} inputs. In other words, the results
  cannot contain any information about the input, which clearly means no utility
  is possible.
\end{itemize}

On the other hand, when $\epsilon$ tends to infinity, then the inequality is
automatically satisfied because the right-hand side become arbitrarily large.
Therefore, no constraint is imposed on the function and privacy drops off
very quickly.  This yields maximal utility, because we can have all available
data without randomness.

\subsection{Interpreting the Definition}

Now that we have thought about the definition and some of its technical
implications, let us think about what it means for privacy.

\paragraph{Inference and protection from harm.}
One view of privacy is that it is about protecting individuals from harm
that may arise from the release of their data. By learning things about
individuals, a party with corrupt intent might use that information to
limit their opportunities (e.g., deny them a job or a loan), offer
differentiated services (e.g., higher prices for customers from affluent
areas), or otherwise discriminate against them in numerous ways that
play against their advantage.

One question that we might ask is, why not strive for a definition that prevents
such parties from learning \emph{anything} new about an individual from a result
involving their data? If nothing new about the individual can be learned from
the release, then no harm can follow. Researchers have contemplated this
possibility before~\citep{Dwork06icalp}, and not suprisingly it turns out that
doing so is at fundamental odds with a simultaneous goal of extracting useful
insights from personal information.

Differential privacy aims to protect individuals from such harm to the greatest
extent possible. The key to this is the \emph{relative} nature of the
definition. Rather than trying to prevent users from learning \emph{anything}
about an individual, we can think of the definition as trying to prevent users
from learning new things about an individual relative to what they \emph{could
  have} learned had the individual not shared their data. This is where the idea
of neighboring inputs comes from: a neighboring input is one in which a
particular individual's data takes a different value, which we can view as being
a input where everyone \emph{except} that individual shared (i.e., some other
individual took their place). Differential privacy requires that any output of
$f$ be approximately as likely in both cases: one where the individual shared
their data, and one where they did not.

For example, suppose that you are given the opportunity to share your medical
records with a researcher who will use them in a study intended to improve
treatments. You may rightly be concerned that if the researcher publishes
results based on your data, a data-savvy insurance provider might be able to
infer something about your health status from these results in the future, and
decide to raise your premiums or deny coverage. However, if the researcher
applied differential privacy with an appropriately-chosen $\epsilon$, then you
might be reassured that no results that could come of the study would be that
much more or less likely because of your decision to share. It follows that if
an insurer were to base their decision on those differentially-private results,
then they are similarly not much more or less likely to deny you coverage.

\paragraph{Plausible deniability.}
Another way of looking at the protection given by differential privacy is in
terms of \emph{plausible deniability}, or one's ability to make a believable
claim that their data takes some value of their choosing, i.e., to ``deny'' a
claim that their data took the value it did. Because \autoref{def:dp} requires
that the likelihood of $f$ responding with any value $s$ is nearly identical
regardless of what value the individual's data took, it would indeed be
reasonable for the individual to claim that their data took another value; the
probability of producing $s$ would be about the same no matter what value they
chose.

\paragraph{Indistinguishability and influence.}
Another way of viewing the definition, which brings us closer to the semantics
of the computation done by $f$, is in terms of how much individuals' data can
influence, or cause changes to, $f$'s response. We have talked about influence
before in the context of noninterference, which required that the $\ms{H}$-typed
parts of the initial state have no influence on the $\ms{L}$-typed parts of the
final state:

\begin{quote}
  For all $\omega_1$, $\omega_2$,
  $\Sigma \vdash \omega_1 \approx_{\ms{L}} \omega_2$ implies
  $\Sigma \vdash \eval\; \omega_1\; \alpha \approx_{\ms{L}} \eval\; \omega_2\;
  \alpha$
\end{quote}
We might rewrite \autoref{def:dp} more concisely as follows.
\begin{quote}
  For all $m_1$, $m_2$, $m_1 \sim_1 m_2$ implies
  $\pr[f(m_1) = s] \leq e^\epsilon \cdot \pr[f(m_2) = s]$
\end{quote}

Notice the similarities between these definitions:

\begin{itemize}
\item In both cases, the definitions quantify over all pairs of inputs
  (i.e., initial states) that are related in a way that reflects what we
  are trying to protect. For noninterference, the relation does this by
  only constraining the low-security variables, so that the final state is
  indistinguishable regardless of the initial high-security variables. For
  differential privacy, the neighbor relation works similarly by letting
  one individual's data take an arbitrary value, and fixing the rest of
  the input.

\item The right-hand side of the implication in each case describes the sort of
  changes that inputs, and more precisely inputs described by the left-hand
  side, are allowed to cause. Noninterference rules out any changes to
  low-security variables, whereas differential privacy places limits on the
  probability of variation in the response.
\end{itemize}

Viewed this way, differential privacy is a property which states that the
influence of individual indices on $f$'s response should remain low, so that
responses computed under neighboring inputs are ``almost''
indistinguishable. This is the essential property that allows for plausible
deniability and protection from harm, and the core of differential privacy's
strong guarantees.

Recall also that we were able to prove that programs satisfy noninterference,
even to the point of designing type systems that simplify the task of writing
noninterferent programs, and can be checked efficiently. Given the similarity
between these definitions, it should not be too surprising that we can also
prove program's adherence to differential privacy. This is part of the appeal of
using the definition in practice: it provides a crisp mathematical formulation
of what it means to be private, that can be proved on real computations.

\subsection{Proving differential privacy: randomized response}

Now let us go back to our example of randomized response. Does it satisfy
differential privacy? Let us keep things simple and assume that $F$ is
the identity function that just returns the contents of $m[0]$, $p
= 1/2$ and all variables and memory cells hold values in the set
$\{0,1\}$. This corresponds to the following program:
\[
  \begin{array}{l}
    f(m) = \\
\ \ b := \flip(p) \\
\ \ \mathbf{if}\ b = 1\ \mathbf{then} \\
\ \ \ \ \mi{out} := m[0] \\
\ \ \mathbf{else} \\
\ \ \ \ \mi{out} := \flip(p)
  \end{array}
\]
It turns out that this does indeed satisfy $\epsilon$-differential privacy.
Normally, we would do our calculation and then find a tight $\epsilon$, but
let's anticipate it.

\begin{theorem}
\label{thm:randresp}
$f$ satisfies $\mathrm{ln}(3)$-differential privacy when $p = 1/2$
and $m[0] \in \{0,1\}$.
\end{theorem}

\begin{proof}
Recall that we need to show that the following inequality holds over all
pairs of neighboring inputs and all outputs $s$:
\[
\pr[f(m_1) = s] \leq e^\epsilon \cdot \pr[f(m_2) = s]
\]
Because this instantiation of randomized response only depends on the
contents of a single memory cell, i.e. $m[0]$, There are two
possible configurations of neighboring inputs: $m_1[0] = 1,
m_2[0] = 0$ and $m_1[0] = 0, m_2[0] = 1$.

\begin{description}
\item[Case:] $m_1[0] = 1, m_2[0] = 0$. The inequality has to be
  satisfied for all $s$, that is, for $s = 1$ and $s = 0$.
  \begin{description}
  \item[Subcase:] $s = 1$.  Then we calculate the left-hand side
    \[
      \begin{array}{ll}
        \pr[f(m_1) = 1]
        &= \mathrm{Pr}[b = 1] + \mathrm{Pr}[b = 0 \land \flip(p) = 1] \\
        &= p + (1-p)p = 3/4
      \end{array}
    \]
    and the right-hand side
    \[
      \begin{array}{ll}
        \pr[f(m_2) = 1]
        &= \mathrm{Pr}[b = 0 \land \flip(p) = 1] \\
        &= (1-p)p = 1/4
      \end{array}
    \]
    For the inequality to hold we have to have
    \[
      \pr[f(m_1) = 1] = 3/4 \leq e^\epsilon\cdot 1/4 = e^\epsilon \cdot \pr[f(m_2) = 1]
    \]
    We see $3 \leq e^\epsilon$ and therefore $\epsilon = \mathrm{ln}(3)$ is a
    tight bound.
  \item[Subcase:] $s = 0$.  Again, we calculate both sides:
    \[
      \begin{array}{ll}
        \pr[f(m_1) = 0]
        &= \Pr[b = 0 \land \flip(p) = 0] \\
        &= (1-p)^2 = 1/4
      \end{array}
    \]
    And for the other side:
    \[
      \begin{array}{ll}
        \pr[f(m_2) = 0]
        &= \pr[b = 1] + \pr[b = 0 \land \flip(p) = 0] \\
        &= p + (1-p)(1-p) = 3/4
      \end{array}
    \]
    Summarizing
    \[
      \begin{array}{ll}
        \pr[f(m_1) = 0] = 1/4 \leq e^\epsilon \cdot 3/4 = e^\epsilon \cdot \pr[f(m_2) = 0]
      \end{array}
    \]
    so $1/3 \leq e^\epsilon$, which is satisfied for any $\epsilon \geq 0$.
  \end{description}
\item[Case:] $m_1[0] = 0, m_2[0] = 1$.  In the case where $p = 1/2$ this turns
  out to be symmetric to the previous case.
  \begin{description}
  \item[Subcase:] $s = 1$.
    \[
      \begin{array}{ll}
        \pr[f(m_1) = 1]
        &= \pr[b = 0 \land \flip(p) = 1] = (1-p)p  = 1/4 \\
        \pr[f(m_2) = 1]
        &= \pr[b = 1] + \pr[b = 0 \land \flip(p) = 1] \\
        &= p + (1-p)p  = 3/4
      \end{array}
    \]
  \item[Subcase:] $s = 0$.
    \[
      \begin{array}{ll}
        \pr[f(m_1) = 0]
        &= \pr[b = 1] + \pr[b = 0 \land \flip(p) = 0] \\
        &= p + (1-p)p  = 3/4 \\
        \pr[f(m_2) = 1]
        &= \mathrm{Pr}[b = 0 \land \flip(p) = 0] \\
        &= (1-p)(1-p)  = 3/4
      \end{array}
    \]
    So indeed the probabilities are simply inverted in this case
  \end{description}
\end{description}
\end{proof}

Now what would we expect for $p = 1/4$?  We go into the branch where we reveal
$m[0]$ less frequently, instead returning a random (if biased) answer instead.
This would indicate that the privacy budget could be less than $\mathrm{ln}(3)$.
Conversely, if $p = 3/4$ we go into the revealing branch with much higher
probably, so we would need a higher privacy budget to be allowed to do that.

If we did our math right, there would be four inequalities that
should be satisfied for general $p$:
\newcommand{\ffrac}[2]{\frac{\displaystyle #1}{\displaystyle #2}}
\[\renewcommand{\arraystretch}{2}
  \begin{array}{lcl}
    \ffrac{p + (1-p)p}{(1-p)p} \leq e^\epsilon \\
    \ffrac{(1-p)(1-p)}{p+(1-p)(1-p)} \leq e^\epsilon \\
    \ffrac{(1-p)p}{p+(1-p)p} \leq e^\epsilon \\
    \ffrac{p+(1-p)p}{(1-p)(1-p)} \leq e^\epsilon
  \end{array}
\]
The middle two fractions are less than 1 and are therefore automatically
satisfied.

Plugging in $p = 1/4$, we get $\epsilon \geq \mathrm{ln}(7/3)$, which is indeed
less than $\mathrm{ln}(3)$.  Plugging in $p = 3/4$, we get
$\epsilon \geq \mathrm{ln}(15)$, which is indeed greater than $\mathrm{ln}(3)$.
One could also imagine inverting the question: given a certain privacy budget,
can we choose a suitable $p$ to make the function comply?

In \href{\courseurl/lectures/23-dpmechanisms.pdf}{Lecture 23} we will see how
differential privacy is achieved for a much broader class of queries than
randomized response, through general-purpose mechanisms for adding calibrated
noise, and how multiple such queries can be composed while tracking their
cumulative privacy cost.

\section{Practice Problems}

\begin{task}\rm
  The course staff wants to release a statistic about the final scores of
  students in 15-316.  Assume that $X$ is an array of length $N$ containing
  student grades, where each $X(i) \in [0, 100]$.

  After a regrettable lapse in judgement, the instructor proposes the following
  scheme $\alpha$:
  \begin{enumerate}
  \item Pick an index $i$ uniformly at random from the set
    $\{j : X(j) \geq 90\}$.  If this set is empty, return $\bot$.
  \item Otherwise, return the pair $(i, X(i))$.
  \end{enumerate}
  Using \autoref{def:dp}, explain why this scheme cannot be
  $\epsilon$-differentially private for \emph{any} $\epsilon > 0$.  In
  particular, exhibit a pair of neighboring datasets and an output value for
  which the inequality must fail.
\end{task}

\begin{task}\rm
  Consider the following variant of randomized response, where the subject
  flips a biased coin to decide whether to report the true answer or its
  \emph{opposite} (rather than a fresh random bit as in the version above):
  \[
    \begin{array}{l}
      b := \flip(0.6) \\
      \mathbf{if}\ b = 1\ \mathbf{then} \\
      \ \ \mi{out} := x \\
      \mathbf{else} \\
      \ \ \mi{out} := 1 - x
    \end{array}
  \]
  where $x \in \{0, 1\}$.
  \begin{enumerate}
  \item Compute $\pr[f(0) = 0]$, $\pr[f(0) = 1]$, $\pr[f(1) = 0]$, and
    $\pr[f(1) = 1]$.
  \item Determine the tight $\epsilon$ for which this function satisfies
    $\epsilon$-differential privacy.
  \item What happens to the tight $\epsilon$ as the bias approaches
    $p = 1/2$?  Explain why this matches your intuition about the
    privacy--utility tradeoff.
  \end{enumerate}
\end{task}

\bibliographystyle{plainnat}
\bibliography{bibliography}

\end{document}
